\documentclass[10pt,a4paper]{amsart}
\usepackage[margin=19mm]{geometry}
\usepackage{amsmath,amssymb,mathtools}
\usepackage[scr=rsfs,cal=euler]{mathalfa}
\usepackage{xfrac}
\usepackage[unicode,hidelinks,bookmarks=false]{hyperref}
\newcommand{\ie}{i.e.\ }
\newcommand{\cf}{cf.\ }
\newcommand{\resp}{respectively\ }
\newcommand{\Subsection}{\subsection}
\providecommand{\nobreakdash}{\nobreak\hbox{-}}
\setlength{\emergencystretch}{2em}
\multlinegap0pt
\usepackage{harpoon}
\usepackage{tikz}
\usepackage{float}
\usepackage{mathrsfs}
\usepackage{verbatim}
\usepackage{enumitem}
\usetikzlibrary{positioning,chains,fit,shapes,calc}  
\newtheorem{theorem}{Theorem}
\newtheorem{lemma}{Lemma}
\newtheorem{definition}{Definition}
\newtheorem{construction}{Construction}
\newtheorem{remark}{Remark}
\newtheorem{corollary}{Corollary}
\newtheorem{property}{Property}
\newtheorem{requirement}{Requirement}
\newcommand{\vect}[1]{\accentset{\rightharpoonup}{#1}}
\newcommand{\harpoon}{\overset{\rightharpoonup}}
\newcommand{\ppp}{\[\begin{aligned}}
\newcommand{\ooo}{\end{aligned}\]}
\newcommand{\err}{\mathrm{err}}
\newcommand{\rad}{\mathrm{rad}}
\newcommand{\val}{\mathrm{val}}
\newcommand{\ind}{\mathrm{ind}}
\newcommand{\mrm}[1]{\mathrm{#1}}
\newcommand{\sgn}{\mathrm{sgn}}
\newcommand{\degu}{\mathrm{deg}_{u_3}}
\newcommand{\iter}{\mathrm{iter}_\phi}
\newcommand{\maxphi}{\mathrm{max}_\phi}
\newcommand{\maxMult}{\mathrm{maxMult}_\phi}
\newcommand{\RPT}{\mathrm{RPT}}
\begin{document}
\title[Multi-point variants of the Newton-Raphson-Simpson method ar]{Multi-point variants of the Newton-Raphson-Simpson method arising from organizing a formal zero according to a function $\phi$}
\author{Mario DeFranco}
\date{}
\maketitle
\noindent\emph{Source and licence.} Multi-point variants of the Newton-Raphson-Simpson method arising from organizing a formal zero according to a function $\phi$. Version arXiv:2608.12384v1 (CC BY 4.0). This material is distributed under the Creative Commons Attribution 4.0 International licence, http://creativecommons.org/licenses/by/4.0/, as recorded in the arXiv OAI licence metadata for this exact version and shown on the abstract page https://arxiv.org/abs/2608.12384v1. Full rights notice: licenses/arxiv-2608.12384-CC-BY-4.0.txt.\par\smallskip
\noindent\textit{Editorial scope.} Counted unit: the original \texttt{theorem} environment \texttt{t main} at source lines 182--187 together with its complete proof at lines 188--230; the delivered document keeps source lines 106--231 so that every referenced label and every citation resolves inside it. Native editable source is the paper's own concatenated TeX; the AMS wrapper is editorial formatting only. No publisher class, upstream script or unrelated artwork is redistributed.
\section{Max-phi Methods}

For an integers $a\leq b$, let $[a,b]$ denote the set 
\[
[a,b]=\{ i \in \mathbb{Z} \colon a\leq i \leq b\}.
\]

\begin{definition} \label{d iterphi}
 Let $\phi$ be a function 
\[
\phi \colon \mathbb{Z}_{\geq 1} \rightarrow  \mathbb{Z}_{\geq 0}. 
\]
For $T \in \mathrm{RootedPlaneTrees}$, define the iteration number $\iter(T)$ of $T$ by  
\[
\iter(T) = \begin{cases} 0 &\text{ if } T=()\\ 
   \maxphi(T)+ \phi(\maxMult(T)) &\text{ if } T= (T_1, \ldots, T_l)
\end{cases}.
\]
where
\[
\maxphi(T) = \max(\{\iter(T_i) \colon 1 \leq i \leq l \})
\]
and 
\[
\maxMult(T) = \# \{i \colon  1 \leq i \leq l  \text{ and } \iter(T_i) = \maxphi(T)\}.
\]
\end{definition} 

\begin{definition} 
For an integer $N \geq 0$, let $Z_N$ denote 

\[
Z_N = \sum_{\substack{ T \in \mathrm{RootedPlaneTrees} \\ \iter(T) \leq N }} w(T) \\ 
\]

\end{definition}

\begin{lemma} \label{l krsum}
For $N\geq 0$, we have
\begin{align}
Z_N = & -\frac{a_0}{a_1}+\sum_{k=0}^{\infty} \sum_{r \in \phi^{-1}([0,k])} (Z_{N-k}-Z_{N-k-1})^r \sum_{l=2}^d { l \choose r} (-\frac{a_l}{a_1})Z_{N-k-1}^{l-r}. 
         \end{align}
\end{lemma}

\begin{proof}
Let $T \in \mathrm{RootedPlaneTrees}$ with 
\begin{equation} \label{ineq}
\iter(T) \leq N.
\end{equation}
 The term $ -\frac{a_0}{a_1}$ corresponds to $w(T)$ for $T=()$. Thus assume $T = (T_1, \ldots, T_l)$ with $l \geq 2$ and \[\maxphi(T)=N-k\] and \[\maxMult(T) = r\] for some $k \geq 0$ and $r \geq 1$. Then inequality \eqref{ineq} implies $\phi(r) \leq k$ i.e. $r \in \phi^{-1}([0,k])$. The factor $(Z_{N-k}-Z_{N-k-1})^r$ corresponds the weights of the $r$ trees $T_i$ with $\iter(T_i) = \maxphi(T)$ and the factor $Z_{N-k-1}^{l-r}$ corresponds to the weights of the remaining $l-i$ trees. There are ${ l \choose r}$ ways to order these trees, and $(-\frac{a_l}{a_1})$ completes the weight of $T$. This completes the proof. 
 \end{proof}

\begin{lemma} \label{l krsum L}
Suppose that the image of $\phi$ is $[0,L]$ for some integer $L \geq 1$. Then
\begin{align}
Z_N = &\sum_{k=0}^{L-1} \sum_{r \in \phi^{-1}([0,k])} (Z_{N-k}-Z_{N-k-1})^r \sum_{l=2}^d { l \choose r} (-\frac{a_l}{a_1})Z_{N-k-1}^{l-r}  \label{krsum L sum}\\
         & + (-\frac{a_0}{a_1})+\sum_{i=2}^d  (-\frac{a_i}{a_1}) Z_{N-L}^i. \label{krsum L f}
\end{align}
\end{lemma}
\begin{proof}
The assumption on $\phi$ means for each $k \geq L$ that
\[
\phi^{-1}([0,k]) =  [1,\infty).
\]
From Lemma \ref{l krsum}, it thus suffices to prove  
\[
\sum_{k=L}^{\infty} \sum_{r =1}^\infty { l \choose r} (Z_{N-k}-Z_{N-k-1})^r  Z_{N-k-1}^{l-r} =   Z_{N-L}^l.
\]
By the binomial theorem, the left side is 
\begin{align*}
\sum_{k=L}^{\infty} (Z_{N-k}-Z_{N-k-1}+ Z_{N-k-1})^l  -   Z_{N-k-1}^l   = &\sum_{k=L}^{\infty} Z_{N-k}^l  -   Z_{N-k-1}^l \\ 
=& Z_{N-L}^l.
\end{align*}
 This completes the proof. 
\end{proof}


\begin{theorem} \label{t main} 
Suppose that the image of $\phi$ is $[0,L]$ for some integer $L \geq 1$ and also that $\phi^{-1}(\{ 0\} )=\{ 1\}$. Then
\begin{align}
Z_N = Z_{N-1} - \frac{f(Z_{N-L})+ \sum_{k=1}^{L-1} \sum_{r \in \phi^{-1}([0,k])} (Z_{N-k} - Z_{N-k-1})^r \frac{f^{(r)}(Z_{N-k-1})}{r!}}{f'(Z_{N-1})}. \label{main}
\end{align}
\end{theorem}

\begin{proof}
 The assumption that $\phi^{-1}(\{ 0\}) = \{ 1\}$ means $1 \in \phi^{-1}([0,k])$ for all $ k \geq 0$. Also 
 \[
 \sum_{l=2}^d { l \choose r} (-\frac{a_l}{a_1})x^{l-r} = 
 \begin{cases} 
 -\frac{f^{(r)}(x)}{r!a_1} \text{ if } r \geq 2 \\ 
 1 -\frac{f'(x)}{a_1} \text{ if } r = 1
 \end{cases}.
 \]
Now in the statement of Lemma \ref{l krsum L} we make the following simplifications. Line \eqref{krsum L f} is equal to 
 \begin{align}
&-\frac{f(Z_{N-L})}{a_1}\\ 
&+ Z_{N-L} \label{tel 1}.
 \end{align}
The $k=0$ contribution at line \eqref{krsum L sum} is equal to 
\[
(Z_{N}-Z_{N-1})(1-\frac{f'(Z_{N-1})}{a_1})
\]
because the only possible $r$ is $r=1$ from the assumption $\phi^{-1}(\{ 0\}) = \{ 1\}$. 
This is equal to 
\begin{align}
Z_N(1-\frac{f'(Z_{N-1})}{a_1}) \label{Z_N term}\\ 
+Z_{N-1}\frac{f'(Z_{N-1})}{a_1} \\
 -Z_{N-1}, \label{tel 2}
\end{align}

From the sum over the remaining $k$, we separate the $r=1$ contribution as 
\[
\sum_{k=1}^{L-1}\sum_{r \in \phi^{-1}([1,k])} \left( -(Z_{N-k} - Z_{N-k-1})^r \frac{f^{(r)}(Z_{N-k-1})}{r!a_1} \right)+  (Z_{N-k} - Z_{N-k-1})(1-\frac{f'(Z_{N-k-1})}{a_1})
\] 
which is equal to 

\begin{align}
&\sum_{k=1}^{L-1}  \sum_{r \in \phi^{-1}([0,k])}-(Z_{N-k} - Z_{N-k-1})^r \frac{f^{(r)}(Z_{N-k-1})}{r!a_1} \\
+ &\sum_{k=1}^{L-1 }  (Z_{N-k} - Z_{N-k-1}). \label{tel 3}
\end{align}
The terms from lines \eqref{tel 1}, \eqref{tel 2}, \eqref{tel 3} sum to 0. 
We move line \eqref{Z_N term} to the left side of the equation, so the resulting left side then is equal to 
\[
Z_N \frac{f'(Z_{N-1})}{a_1}.
\]
Then divide by $\frac{f'(Z_{N-1})}{a_1}$. The $a_1$ factors cancel out. This completes the proof. 
\end{proof}


\end{document}
